% ============================================================================
% Fichier  : partie1/P01_C03_alignement_valeur_decision.tex
% Titre    : Théories et Modèles : Alignement, Valeur et Décision
% Auteur   : MINKA MI NGUIDJOÏ Thierry Emmanuel
% Version  : 1.0 -- octobre 2026
% ============================================================================
% Chapitre associé au laboratoire GCI-LAB-01 (arbitrage du portefeuille de
% l'ANPP, opération MEFOU).
% ============================================================================

\chapter{Théories et Modèles : Alignement, Valeur et Décision}
\label{chap:alignement}
\chaptermark{Alignement, valeur et décision}

\epigraph{%
    « On voit l'ère de l'ordinateur partout, sauf dans les statistiques
    de productivité. »%
}{Robert \textsc{Solow}, prix Nobel d'économie, 1987 (notre traduction)}

\niveauChapitre{Fondamental}{%
    Maîtrise du Chapitre~\ref{chap:gouv-entreprise} (définition de travail
    de la gouvernance des SI et ses trois finalités).%
}

\begin{boxobjectifs}
\begin{itemize}
    \item Appliquer le modèle d'alignement stratégique de Henderson et
          Venkatraman.
    \item Classer les droits de décision selon les archétypes de Weill et
          Ross.
    \item Expliquer les trois finalités~: réalisation des bénéfices,
          optimisation des risques, optimisation des ressources.
    \item Évaluer un investissement informatique (coût total de possession,
          valeur actuelle nette, score multicritère).
    \item Arbitrer un portefeuille de projets.
\end{itemize}
\end{boxobjectifs}

\bigskip

% ============================================================================
\section*{Question directrice : 14 milliards pour quoi faire~?}
% ============================================================================

À l'ANPP (opération MEFOU, Chapitre~\ref{chap:gouv-entreprise}), le
schéma directeur adopté en 2023 prévoyait neuf projets pour 14~milliards
FCFA. À mi-parcours, trois projets ont consommé 71~\% du budget sans être
achevés. Le nouveau président du conseil d'administration pose trois
questions simples~: « Ces projets servaient-ils la mission de l'agence~?
Qui les a choisis, et selon quels critères~? Combien rapportent-ils, ou
combien de risques évitent-ils~? »

Ces trois questions sont celles de ce chapitre. La première relève de
l'\textbf{alignement stratégique} (section~\ref{sec:c03-alignement}). La
deuxième relève des \textbf{droits de décision}
(section~\ref{sec:c03-weill-ross}). La troisième relève de la
\textbf{valeur}, que le chapitre définit (section~\ref{sec:c03-valeur}),
mesure (section~\ref{sec:c03-evaluation}) et arbitre à l'échelle d'un
portefeuille (section~\ref{sec:c03-portefeuille}). Le laboratoire GCI-LAB-01
de NDAP SUKULU applique la démarche au portefeuille complet de l'ANPP.

\separateur

% ============================================================================
\section{Le Paradoxe de la Productivité}
\label{sec:c03-paradoxe}
% ============================================================================

La phrase de Solow, en épigraphe, a ouvert un long débat. Dans les années
1980, les entreprises américaines investissent massivement dans
l'informatique, sans que la productivité mesurée progresse en proportion.
Les travaux ultérieurs ont largement nuancé ce constat~: les gains
existent, mais ils sont inégaux, différés, et dépendent moins de la
technologie elle-même que de la manière dont l'organisation la met en
œuvre.

Pour la gouvernance, la leçon est décisive~: \textbf{la valeur n'est pas
dans la technologie, elle est dans son usage}. Deux organisations qui
achètent le même logiciel n'en tirent pas la même valeur. Ce qui les
distingue, c'est la manière dont elles choisissent leurs projets,
transforment leurs processus et suivent les bénéfices. C'est précisément
le champ de la gouvernance des SI.

\begin{boxinfo}
\textbf{Les quatre questions de Thorp.} John Thorp a résumé la
gouvernance de la valeur en quatre questions~\cite{thorp1998paradox},
reprises plus tard par le référentiel Val~IT de l'ISACA~:
\emph{fait-on les bonnes choses~?} (la stratégie), \emph{les fait-on de la
bonne manière~?} (l'architecture), \emph{les mène-t-on à bien~?} (la
livraison), \emph{en tire-t-on les bénéfices~?} (la valeur). Les deux
premières relèvent de la gouvernance, les deux dernières du management,
sous le contrôle de la gouvernance.
\end{boxinfo}

% ============================================================================
\section{Le Modèle d'Alignement Stratégique}
\label{sec:c03-alignement}
% ============================================================================

\subsection{Quatre domaines, deux dimensions}

Le modèle d'alignement stratégique (\textit{Strategic Alignment Model},
\sigle{SAM}) de Henderson et Venkatraman~\cite{henderson1993alignment}
reste le cadre de référence de l'alignement. Il distingue quatre domaines,
organisés selon deux dimensions.

\begin{itemize}
    \item La dimension de l'\textbf{ajustement stratégique}
          (\textit{strategic fit}) relie le niveau externe (la stratégie~:
          sur quel marché, avec quelles compétences distinctives, avec
          quelle gouvernance) et le niveau interne (l'infrastructure et les
          processus qui la mettent en œuvre).
    \item La dimension de l'\textbf{intégration fonctionnelle}
          (\textit{functional integration}) relie le domaine de l'activité
          et le domaine des technologies de l'information.
\end{itemize}

\begin{center}
\begin{tikzpicture}[
    dom/.style = {
        draw=CouleurBordDef, fill=CouleurDef, rounded corners=4pt,
        text width=4.2cm, minimum height=1.3cm, align=center,
        font=\footnotesize\bfseries\color{CouleurPrimaire}
    },
    lien/.style = {{Stealth[length=5pt]}-{Stealth[length=5pt]}, thick,
        color=CouleurSecondaire}
]
\node[dom] (bs) at (0,2.6)   {Stratégie de l'organisation};
\node[dom] (ts) at (7.6,2.6) {Stratégie des technologies\\ de l'information};
\node[dom] (oi) at (0,0)     {Infrastructure et processus\\ de l'organisation};
\node[dom] (ii) at (7.6,0)   {Infrastructure et processus\\ des SI};
\draw[lien] (bs) -- node[above, font=\scriptsize\itshape, align=center]{intégration\\
    fonctionnelle} (ts);
\draw[lien] (oi) -- node[below, font=\scriptsize\itshape, align=center]{intégration\\
    fonctionnelle} (ii);
\draw[lien] (bs) -- node[left, font=\scriptsize\itshape, align=center]{ajustement\\
    stratégique} (oi);
\draw[lien] (ts) -- node[right, font=\scriptsize\itshape, align=center]{ajustement\\
    stratégique} (ii);
\node[font=\scriptsize\bfseries\color{CouleurAccent}, above=0.1cm of bs]
    {Activité};
\node[font=\scriptsize\bfseries\color{CouleurAccent}, above=0.1cm of ts]
    {Technologies};
\end{tikzpicture}
\end{center}
\vspace{-0.4em}
\begin{center}{\small\itshape Figure~3.1. Le modèle d'alignement
stratégique (d'après Henderson et Venkatraman, 1993)}\end{center}

\begin{boxdef}
\termecle{Alignement stratégique}\index{alignement stratégique}
(\textit{strategic alignment})~: degré de cohérence entre la stratégie de
l'organisation, la stratégie des technologies de l'information et les
infrastructures et processus qui les mettent en œuvre. Pour Henderson et
Venkatraman, l'alignement n'est pas un état atteint une fois pour toutes,
mais un processus continu d'adaptation.
\end{boxdef}

\subsection{Quatre perspectives d'alignement}

Le modèle décrit quatre manières de parcourir les domaines, selon le point
de départ et le point d'arrivée. Chacune correspond à un rôle différent
de la direction générale et du DSI.

\begin{tableaucours}{L{3.0cm}L{4.4cm}L{4.8cm}}{%
    \tableauHeader{Perspective} &
    \tableauHeader{Parcours} &
    \tableauHeader{Exemple camerounais}
}
Exécution de la stratégie & Stratégie de l'organisation $\rightarrow$
infrastructure de l'organisation $\rightarrow$ infrastructure des SI &
L'ANPP doit payer tous les fournisseurs de l'État sous 30~jours~: elle
réorganise ses circuits, puis outille la chaîne de paiement \\
Transformation technologique & Stratégie de l'organisation $\rightarrow$
stratégie des TI $\rightarrow$ infrastructure des SI & La BANQUE SANAGA
veut conquérir les jeunes urbains~: elle choisit une stratégie mobile,
puis la plateforme qui la porte \\
Potentiel concurrentiel & Stratégie des TI $\rightarrow$ stratégie de
l'organisation $\rightarrow$ infrastructure de l'organisation & NOUN PAY
découvre que la notation par IA permet le microcrédit instantané~: elle
redéfinit son offre \\
Niveau de service & Stratégie des TI $\rightarrow$ infrastructure des SI
$\rightarrow$ infrastructure de l'organisation & KRIBILOG fiabilise son
suivi des conteneurs, ce qui permet aux équipes du port de réorganiser
leurs rotations \\
\end{tableaucours}
\vspace{-0.8em}
\begin{center}{\small\itshape Tableau~3.1. Les quatre perspectives
d'alignement}\end{center}

\begin{boxalert}
\textbf{Violation typique.} Un projet est lancé parce que « la
technologie existe » ou parce qu'un autre organisme l'a adoptée, sans que
personne n'ait identifié l'objectif de l'organisation qu'il sert. Le
projet est livré, mais aucun processus ne change, et aucun bénéfice n'est
mesuré.

\smallskip
\textbf{Règle~:} tout projet doit pouvoir être rattaché, par écrit, à un
objectif de l'organisation et à l'une des quatre perspectives. Un projet
qu'on ne peut rattacher à rien n'est pas aligné~: il doit être justifié
autrement, ou abandonné.
\end{boxalert}

% ============================================================================
\section{Qui Décide de Quoi : les Archétypes de Weill et Ross}
\label{sec:c03-weill-ross}
% ============================================================================

\subsection{Cinq décisions}

Weill et Ross~\cite{weill2004itgovernance} identifient cinq grandes
catégories de décisions relatives aux technologies de l'information.
Elles sont liées~: les principes encadrent les quatre autres.

\begin{enumerate}
    \item \textbf{Les principes}~: le rôle que les TI doivent jouer dans
          l'organisation (réduire les coûts, soutenir la croissance,
          transformer le métier).
    \item \textbf{L'architecture}~: l'organisation logique des données, des
          applications et de l'infrastructure, et les standards retenus.
    \item \textbf{L'infrastructure}~: les services techniques partagés
          (réseaux, centres de données, sécurité) et leur niveau.
    \item \textbf{Les besoins applicatifs}~: les applications à acquérir ou
          à développer pour les métiers.
    \item \textbf{L'investissement et la priorisation}~: combien investir,
          dans quoi, et dans quel ordre.
\end{enumerate}

\subsection{Six archétypes}

Pour chaque décision, Weill et Ross se demandent \emph{qui} détient le
droit de décider, et \emph{qui} est consulté avant. Ils distinguent six
archétypes de gouvernance.

\begin{tableaucours}{L{2.8cm}L{4.6cm}L{4.8cm}}{%
    \tableauHeader{Archétype} &
    \tableauHeader{Qui décide} &
    \tableauHeader{Risque typique}
}
Monarchie métier & La direction générale ou un groupe de dirigeants des
métiers & Décisions déconnectées des contraintes techniques \\
Monarchie informatique & Les responsables informatiques & Décisions
techniquement justes, mais sans adhésion des métiers \\
Féodalité & Chaque direction métier pour son propre périmètre & Doublons,
incompatibilités, coûts non mutualisés \\
Fédéralisme & La direction générale et les métiers ensemble, avec ou sans
l'informatique & Lenteur, compromis difficiles \\
Duopole & L'informatique et un groupe de métiers & Dépend de la qualité
de la relation bilatérale \\
Anarchie & Chaque utilisateur ou petit groupe & Aucune cohérence, risques
non maîtrisés \\
\end{tableaucours}
\vspace{-0.8em}
\begin{center}{\small\itshape Tableau~3.2. Les six archétypes de Weill et
Ross}\end{center}

\subsection{La matrice des droits de décision}

L'outil le plus utile est la \termecle{matrice des droits de
décision}\index{matrice des droits de décision}~: en lignes les
archétypes, en colonnes les cinq décisions, et dans chaque case qui
apporte un avis (\emph{consultation}) et qui tranche (\emph{décision}).
Elle décrit la gouvernance telle qu'elle fonctionne réellement, et non
telle qu'elle est écrite.

\begin{boxscenario}{Opération MEFOU. Les droits de décision réels}
L'entretien avec les responsables de l'ANPP fait apparaître la situation
suivante.

\medskip
\begin{tableaucours}{L{3.6cm}L{4.0cm}L{4.4cm}}{%
    \tableauHeader{Décision} &
    \tableauHeader{Qui décide réellement} &
    \tableauHeader{Archétype}
}
Principes & Personne (aucun document) & Absence de décision \\
Architecture & Chaque prestataire, projet par projet & Anarchie \\
Infrastructure & Cheffe de la division informatique & Monarchie
informatique \\
Besoins applicatifs & Chaque direction, avec son prestataire & Féodalité \\
Investissement & Directeur général, à la demande de la tutelle &
Monarchie métier \\
\end{tableaucours}

\medskip
\textbf{Analyse} (confiance élevée). La décision la plus lourde,
l'investissement, est prise par une seule personne, sans consultation de
la division informatique ni critères écrits. Les besoins applicatifs sont
définis direction par direction, ce qui explique que trois projets aient
absorbé l'essentiel du budget sans cohérence d'ensemble. L'absence de
principes rend tout arbitrage impossible~: il n'existe aucun critère pour
dire non.
\end{boxscenario}

\begin{boiteRemarque}{Existe-t-il un bon archétype~?}
Il n'existe pas d'archétype bon en soi. Weill et Ross observent que les
organisations les plus performantes combinent des archétypes différents
selon les décisions~: par exemple une monarchie informatique pour
l'architecture et l'infrastructure, où l'expertise prime, et un
fédéralisme ou un duopole pour les besoins applicatifs et
l'investissement, où l'adhésion des métiers est indispensable. La seule
configuration à proscrire est celle qui n'a pas été choisie~: la
gouvernance de fait, qui s'installe faute de décision.
\end{boiteRemarque}

% ============================================================================
\section{Les Trois Finalités de la Gouvernance}
\label{sec:c03-valeur}
% ============================================================================

\subsection{Créer de la valeur}

COBIT~2019 définit l'objectif de la gouvernance comme la \textbf{création
de valeur}, qui résulte de l'équilibre entre trois
finalités~\cite{isaca2018cobitintro}. La définition de travail du
Chapitre~\ref{chap:gouv-entreprise} les reprend.

\begin{boxdef}
\termecle{Réalisation des bénéfices}\index{réalisation des bénéfices}
(\textit{benefits realization})~: créer de la valeur pour l'organisation
grâce aux technologies de l'information, en entretenant et en accroissant
la valeur des investissements existants, et en éliminant les initiatives
et les actifs qui n'en créent pas assez.

\medskip
\termecle{Optimisation des risques}\index{optimisation des risques}
(\textit{risk optimization})~: traiter les risques liés à l'usage, à la
propriété, à l'exploitation et à l'adoption des technologies de
l'information, en les maintenant dans les limites de l'appétence au
risque fixée par l'organe de gouvernance.

\medskip
\termecle{Optimisation des ressources}\index{optimisation des
ressources} (\textit{resource optimization})~: s'assurer que les
compétences et les ressources nécessaires (personnes, processus,
technologies, données) sont disponibles, suffisantes et utilisées de
manière efficace.
\end{boxdef}

Le mot clé est \textbf{équilibre}. Maximiser les bénéfices sans se soucier
des risques conduit à l'opération SANAGA. Minimiser les risques sans se
soucier des bénéfices conduit à l'immobilisme. Économiser les ressources
sans se soucier des deux autres conduit à des systèmes vieillissants et
fragiles.

\subsection{Trois objectifs de gouvernance}

À chaque finalité correspond un objectif du domaine de gouvernance de
COBIT~2019. Ils seront étudiés au Chapitre~\ref{chap:cobit-objectifs}.

\begin{tableaucours}{L{3.4cm}L{4.4cm}L{4.4cm}}{%
    \tableauHeader{Finalité} &
    \tableauHeader{Objectif COBIT 2019} &
    \tableauHeader{Question au conseil}
}
Réalisation des bénéfices & EDM02 · Assurer la réalisation des bénéfices
& Les investissements produisent-ils les bénéfices annoncés~? \\
Optimisation des risques & EDM03 · Assurer l'optimisation des risques &
Les risques sont-ils dans les limites que nous avons fixées~? \\
Optimisation des ressources & EDM04 · Assurer l'optimisation des
ressources & Avons-nous les compétences et les moyens de notre
stratégie~? \\
\end{tableaucours}
\vspace{-0.8em}
\begin{center}{\small\itshape Tableau~3.3. Finalités et objectifs de
gouvernance (intitulés traduits par nos soins)}\end{center}

\begin{boxobjet}{EDM02 · Assurer la réalisation des bénéfices}
\textbf{Finalité~:} obtenir la valeur optimale des initiatives, des
services et des actifs informatiques, et une gestion des coûts efficace.

\textbf{Pratiques~:} évaluer l'optimisation de la valeur~; diriger
l'optimisation de la valeur (critères d'investissement, rôles)~;
surveiller l'optimisation de la valeur (indicateurs, écarts).

\textbf{Sorties~:} principes et critères d'investissement, retour sur
le portefeuille, actions correctives.

\textbf{RACI (extrait)~:} \emph{A} au conseil d'administration, \emph{R}
au directeur général et au comité stratégique des SI, \emph{C} au
directeur financier et au DSI.

\textbf{Indicateurs~:} part des investissements dont les bénéfices sont
suivis après la mise en service~; écart entre bénéfices annoncés et
bénéfices constatés.
\end{boxobjet}

% ============================================================================
\section{Évaluer un Investissement Informatique}
\label{sec:c03-evaluation}
% ============================================================================

\subsection{Le coût total de possession}

Le prix d'achat d'un logiciel n'est qu'une partie de son coût. Le
\termecle{coût total de possession}\index{coût total de possession}
(\textit{total cost of ownership}, \sigle{TCO}) additionne tous les coûts
supportés sur la durée de vie de la solution~: acquisition, intégration,
matériel, formation, exploitation, support et sortie.

\begin{tableaucours}{L{8.6cm}r}{%
    \tableauHeader{Poste (sur 5 ans)} &
    \tableauHeader{Millions FCFA}
}
Licences (prix affiché par l'éditeur) & 300 \\
Intégration et reprise des données & 250 \\
Matériel et hébergement & 150 \\
Formation des utilisateurs & 50 \\
Exploitation interne (90 par an) & 450 \\
Support et maintenance de l'éditeur (60 par an) & 300 \\
\textbf{Coût total de possession} & \textbf{1~500} \\
\end{tableaucours}
\vspace{-0.8em}
\begin{center}{\small\itshape Tableau~3.4. Coût total de possession d'un
logiciel de gestion des paiements (exemple)}\end{center}

Le coût total est cinq fois le prix affiché. C'est un ordre de grandeur
fréquent, et l'une des raisons pour lesquelles les budgets des schémas
directeurs dérapent~: ils ont été construits sur les prix d'acquisition.

\subsection{La valeur actuelle nette}

Un franc reçu dans quatre ans vaut moins qu'un franc dépensé aujourd'hui.
La \termecle{valeur actuelle nette}\index{valeur actuelle nette}
(\sigle{VAN}, \textit{net present value}) ramène tous les flux d'un projet
à leur valeur d'aujourd'hui, au moyen d'un taux d'actualisation $r$~:
\[
    \mathrm{VAN} = -I_0 + \sum_{t=1}^{n} \frac{F_t}{(1+r)^t}
\]
où $I_0$ est l'investissement initial et $F_t$ le flux net (bénéfices
moins coûts d'exploitation) de l'année~$t$. Un projet dont la VAN est
positive crée de la valeur au taux retenu.

Lorsque le flux annuel $F$ est constant, la somme se calcule avec le
facteur d'annuité~:
\[
    \sum_{t=1}^{n} \frac{1}{(1+r)^t} = \frac{1-(1+r)^{-n}}{r},
    \qquad \text{soit } 3{,}7908 \text{ pour } r=10~\% \text{ et } n=5.
\]

\begin{boxscenario}{Opération MEFOU. Trois projets à évaluer}
L'ANPP retient un taux d'actualisation de 10~\% et un horizon de cinq ans.
Trois projets sont en concurrence.

\medskip
\begin{tableaucours}{L{4.4cm}rrr}{%
    \tableauHeader{Projet} &
    \tableauHeader{$I_0$} &
    \tableauHeader{$F$ annuel} &
    \tableauHeader{VAN}
}
A. Paiement dématérialisé des fournisseurs & 1~200 & 450 & $+506$ \\
B. Refonte du portail des ressources humaines & 800 & 180 & $-118$ \\
C. Site de secours (continuité) & 600 & 120 & $-145$ \\
\end{tableaucours}
{\small Montants en millions FCFA.}

\medskip
\textbf{Calcul du projet A.} $450 \times 3{,}7908 = 1~705{,}9$~; VAN $=
1~705{,}9 - 1~200 = 505{,}9$, soit environ $+506$~millions FCFA. Le délai
de récupération est de $1~200 / 450 \approx 2{,}7$~ans, et le taux de
rentabilité interne, pour lequel la VAN s'annule, est d'environ 25~\%.

\medskip
\textbf{Le cas du projet C.} Le site de secours ne produit aucune recette.
Son flux annuel est une \emph{perte évitée}~: l'ANPP estime à 5~\% par an
la probabilité d'un sinistre majeur sur son centre de données, pour une
perte de 2~400~millions FCFA (pénalités de retard de paiement, reprise
manuelle), soit une perte annuelle attendue de $0{,}05 \times 2~400 =
120$~millions FCFA.
\end{boxscenario}

\subsection{Les limites de l'évaluation financière}

Si l'ANPP ne retenait que la VAN, elle financerait le seul projet~A. Ce
serait une erreur de gouvernance, pour trois raisons.
\begin{itemize}
    \item \textbf{Les risques se chiffrent mal.} La perte attendue du
          projet~C repose sur une probabilité très incertaine. Une
          variation de quelques points change le signe de la VAN.
    \item \textbf{Certains projets sont obligatoires.} Une obligation
          légale ou réglementaire ne se discute pas en VAN~: elle se
          respecte, au moindre coût.
    \item \textbf{L'alignement ne se chiffre pas.} Un projet peut avoir
          une VAN modeste et être indispensable à la mission de
          l'organisation.
\end{itemize}

\subsection{Le score multicritère}

Pour dépasser ces limites, l'organe de gouvernance fixe à l'avance des
critères pondérés, puis note chaque projet. Les poids traduisent ses
priorités~: les fixer est une décision de gouvernance, qui ne doit pas
être laissée à ceux qui proposent les projets.

\begin{tableaucours}{L{3.8cm}crrr}{%
    \tableauHeader{Critère} &
    \tableauHeader{Poids} &
    \tableauHeader{A} &
    \tableauHeader{B} &
    \tableauHeader{C}
}
Alignement sur la mission & 30~\% & 5 & 3 & 3 \\
Valeur financière (VAN) & 25~\% & 5 & 2 & 1 \\
Réduction des risques & 20~\% & 2 & 1 & 5 \\
Obligation réglementaire & 15~\% & 3 & 1 & 4 \\
Faisabilité & 10~\% & 4 & 5 & 3 \\
\textbf{Score pondéré} & & \textbf{4,00} & \textbf{2,25} & \textbf{3,05} \\
\end{tableaucours}
\vspace{-0.8em}
\begin{center}{\small\itshape Tableau~3.5. Score multicritère des trois
projets de l'ANPP (notes de 1 à 5)}\end{center}

Le calcul du projet~A est le suivant~: $0{,}30 \times 5 + 0{,}25 \times 5 +
0{,}20 \times 2 + 0{,}15 \times 3 + 0{,}10 \times 4 = 4{,}00$. Avec une
enveloppe de 1~900~millions FCFA, l'ANPP finance A et C (1~800~millions)
et écarte B. Le site de secours, rejeté par la VAN seule, est retenu au
titre des risques et de l'obligation de continuité du service public.

\begin{boxprocedure}
\textbf{Une décision d'investissement auditable.} Pour qu'un auditeur, ou
la Chambre des comptes, puisse établir que le portefeuille a été arbitré,
le dossier doit contenir~:
\begin{itemize}
    \item les critères et leurs poids, approuvés par l'organe de
          gouvernance \emph{avant} l'examen des projets~;
    \item pour chaque projet, un dossier d'investissement (objectif servi,
          coût total de possession, bénéfices attendus et leur
          responsable, risques)~;
    \item la grille de notation, signée par ceux qui l'ont établie~;
    \item le procès-verbal de la décision, y compris pour les projets
          écartés~;
    \item le plan de suivi des bénéfices après la mise en service.
\end{itemize}
\end{boxprocedure}

% ============================================================================
\section{Arbitrer un Portefeuille}
\label{sec:c03-portefeuille}
% ============================================================================

\subsection{Projet, programme, portefeuille}

\begin{boxdef}
\termecle{Portefeuille}\index{portefeuille} (\textit{portfolio})~:
ensemble des projets, des programmes et des services informatiques d'une
organisation, géré comme un tout pour maximiser la valeur obtenue sous
contrainte de ressources et de risque. L'unité de décision de la
gouvernance n'est pas le projet isolé, mais le portefeuille.
\end{boxdef}

Raisonner en portefeuille change la question posée. Il ne s'agit plus de
savoir si un projet est bon, mais s'il est \emph{meilleur que les autres
usages possibles des mêmes ressources}. C'est ce qui a manqué à l'ANPP~:
chaque projet a été approuvé seul, sans jamais être comparé aux autres.

\subsection{Équilibrer le portefeuille}

Weill et Broadbent proposent de classer les investissements informatiques
selon leur finalité~\cite{weill1998infrastructure}. Un portefeuille
équilibré combine les quatre catégories, dans des proportions qui
reflètent la stratégie.

\begin{tableaucours}{L{2.8cm}L{4.6cm}L{4.8cm}}{%
    \tableauHeader{Catégorie} &
    \tableauHeader{Finalité} &
    \tableauHeader{Exemple à l'ANPP}
}
Infrastructure & Fournir les services partagés sur lesquels tout repose &
Réseau sécurisé entre l'agence et les trésoreries régionales \\
Transactionnel & Réduire les coûts et automatiser les opérations
répétitives & Paiement dématérialisé des fournisseurs (projet A) \\
Informationnel & Mieux piloter, mieux contrôler & Tableau de bord des
délais de paiement \\
Stratégique & Changer la position ou la mission de l'organisation &
Paiement des usagers par téléphone mobile \\
\end{tableaucours}
\vspace{-0.8em}
\begin{center}{\small\itshape Tableau~3.6. Les quatre catégories
d'investissement (d'après Weill et Broadbent)}\end{center}

\subsection{Le cadre de la gestion publique au Cameroun}

Pour une administration ou un établissement public, l'arbitrage du
portefeuille s'inscrit dans le cadre budgétaire de l'État.

\begin{boxjuris}
\textbf{\loicam{n°~2018/012} du 11~juillet~2018 portant régime financier
de l'État et des autres entités publiques, article~69.}

\medskip
Selon la présentation qu'en fait le ministère des Finances~\cite{minfi2023regimefinancier},
la loi organise le budget par programmes et institue un
\emph{responsable de programme}, désigné par le ministre. Ce responsable
formule le programme dans le \emph{Projet de performance annuel}
(\sigle{PPA}), affecte les ressources en son sein, organise le dialogue
de gestion et élabore le \emph{Rapport annuel de performance}
(\sigle{RAP}). Il veille au respect des dispositifs de contrôle interne et
de contrôle de gestion.

\medskip
\textbf{Portée.} La gestion publique camerounaise est désormais pilotée
par la performance~: objectifs, indicateurs, compte rendu annuel. Un
projet informatique public doit donc pouvoir être rattaché à un programme
et à ses indicateurs de performance. C'est l'équivalent public de
l'alignement stratégique, et un argument fort pour imposer, à l'ANPP
comme ailleurs, des dossiers d'investissement et un suivi des bénéfices.
\end{boxjuris}
% A VERIFIER au Journal officiel : libellé exact de l'article 69 de la loi
% n° 2018/012 (source utilisée : publication du MINFI, collection
% « Les essentiels de la réforme »).

\begin{boxmetier}{DSI}
\textbf{Présenter un portefeuille au comité stratégique}
\begin{enumerate}
    \item Une page par projet~: objectif servi, perspective d'alignement,
          coût total de possession, VAN ou perte évitée, risques, responsable
          des bénéfices.
    \item Une vue d'ensemble~: répartition du budget entre les quatre
          catégories (Tableau~3.6) et comparaison avec la stratégie.
    \item Le résultat de la grille multicritère, avec les poids approuvés
          par le comité.
    \item Une recommandation explicite~: financer, différer, arrêter.
\end{enumerate}
Le dernier point est le plus difficile, et le plus important~: proposer
d'arrêter un projet en cours est souvent la meilleure décision de
gouvernance qu'un DSI puisse provoquer.
\end{boxmetier}

\begin{boxmetier}{Auditeur SI}
\textbf{Vérifier qu'un portefeuille est réellement arbitré}
\begin{enumerate}
    \item Demander les critères d'investissement, et vérifier leur date
          d'approbation au regard des dates de décision des projets.
    \item Prendre un échantillon de projets approuvés \emph{et} de projets
          écartés~: une organisation qui n'a jamais écarté un projet
          n'arbitre pas.
    \item Pour les projets en service, demander le suivi des bénéfices
          annoncés dans le dossier d'investissement.
\end{enumerate}
\end{boxmetier}

% ============================================================================
\separateur
% ============================================================================

\begin{boxresume}
\textbf{Ce qu'il faut retenir de ce chapitre~:}
\begin{itemize}
    \item \textbf{La valeur est dans l'usage}~: le paradoxe de la
          productivité montre que la technologie seule ne crée pas de
          valeur~; la manière de la choisir et de l'exploiter, si.
    \item \textbf{L'alignement stratégique} relie quatre domaines (stratégie
          et infrastructure, de l'activité et des TI). Il se parcourt selon
          quatre perspectives, et c'est un processus continu.
    \item \textbf{Les droits de décision} portent sur cinq décisions
          (principes, architecture, infrastructure, besoins applicatifs,
          investissement)~; six archétypes décrivent qui décide. Le pire
          archétype est celui qu'on n'a pas choisi.
    \item \textbf{Trois finalités en équilibre}~: réalisation des
          bénéfices, optimisation des risques, optimisation des ressources
          (EDM02, EDM03, EDM04).
    \item \textbf{Évaluer}~: coût total de possession, VAN, délai de
          récupération~; mais aussi perte évitée pour les projets de
          risque, et score multicritère aux poids fixés par la gouvernance.
    \item \textbf{Arbitrer un portefeuille}~: comparer les projets entre
          eux, équilibrer les catégories, et savoir arrêter un projet.
\end{itemize}
\end{boxresume}

\bigskip

\begin{boxexo}[title={\ding{45}\quad Exercice~3.1. Calculer une VAN \niveaubase}]
KRIBILOG SA envisage un logiciel de planification des quais~:
investissement de 450~millions FCFA, gain net de 140~millions FCFA par an
pendant cinq ans, taux d'actualisation de 12~\%.
\begin{enumerate}[(a)]
    \item Calculez le facteur d'annuité pour $r=12~\%$ et $n=5$.
    \item Calculez la VAN du projet. Est-il créateur de valeur~?
    \item Calculez son délai de récupération simple.
\end{enumerate}
\end{boxexo}

\begin{boxexo}[title={\ding{45}\quad Exercice~3.2. La matrice de la BANQUE SANAGA \niveaumoyen}]
À partir des faits du prologue (Chapitre~\ref{chap:prologue}), construisez
la matrice des droits de décision de la BANQUE SANAGA pour le projet
HORIZON.
\begin{enumerate}[(a)]
    \item Pour chacune des cinq décisions de Weill et Ross, indiquez qui
          a réellement décidé et l'archétype correspondant.
    \item Proposez la matrice cible, en justifiant l'archétype retenu pour
          chaque décision.
    \item Rédigez la constatation et l'analyse correspondantes, avec leur
          niveau de confiance.
\end{enumerate}
\end{boxexo}

\begin{boxexo}[title={\ding{45}\quad Exercice~3.3. Arbitrer sous contrainte \niveauexpert}]
Le comité stratégique de l'ANPP révise ses poids~: l'obligation
réglementaire passe à 25~\% et la valeur financière à 15~\%, les autres
poids restant inchangés. Un quatrième projet~D (tableau de bord des
délais de paiement, 300~millions FCFA) obtient les notes 4, 2, 2, 5, 5.
\begin{enumerate}[(a)]
    \item Recalculez les scores des projets A, B, C et D.
    \item Avec une enveloppe de 1~900~millions FCFA, quel portefeuille
          recommandez-vous~? Justifiez.
    \item Le directeur général conteste le changement des poids, qui
          « favorise » le projet~D. Qui doit trancher, et comment
          l'auditeur saura-t-il que les poids n'ont pas été fixés après
          coup pour justifier un choix déjà fait~?
\end{enumerate}
\end{boxexo}

\bigskip

\begin{boxinfo}
\textbf{Transition.} Ce chapitre a posé les modèles qui guident les
décisions de gouvernance~: alignement, droits de décision, valeur. Pour
les mettre en œuvre, les organisations s'appuient sur des référentiels.
Le Chapitre~\ref{chap:referentiels} dresse le panorama de ces
référentiels (COBIT, ITIL, ISO, COSO, TOGAF\ldots) et montre comment ils
s'articulent.
\end{boxinfo}
